\section{\texttt{vmm\_init}：初始化全流程}

\texttt{vmm\_init()} 负责用正式的 4\,KiB 页表替换 \texttt{boot.asm} 的临时 4\,MiB PSE 映射。

\subsection{进入时的 CPU 状态}

\begin{cpustate}
\begin{itemize}
  \item \reg{CR0}.PG = 1（\texttt{boot.asm} 已开启分页，使用 4\,MiB PSE 大页）
  \item \reg{CR3} 指向 \texttt{boot\_pd}（临时页目录）
  \item EIP 在虚拟地址 \hex{C01xxxxx}（high-half）
  \item PMM 已初始化完毕，可以调用 \texttt{pmm\_alloc\_page()}
\end{itemize}
\end{cpustate}

\subsection{核心步骤}

\begin{enumerate}
  \item \textbf{清零页目录}：\texttt{memzero\_page(\&g\_kernel\_pd)} \\
        所有 1024 个 PDE 初始化为 0（Present=0）。

  \item \textbf{计算映射范围}：
        取内核物理结束地址、PMM 管理范围上限，至少 16\,MiB，上限 1\,GiB。

  \item \textbf{建立双重映射}：对 $[0, \text{map\_end})$ 的每个 4\,KiB 页：
    \begin{itemize}
      \item Identity mapping: \texttt{vmm\_map\_page(addr, addr, ...)}
      \item High-half mapping: \texttt{vmm\_map\_page(addr + \hex{C0000000}, addr, ...)}
    \end{itemize}

  \item \textbf{切换 CR3}：
        \texttt{vmm\_load\_page\_directory(VIRT\_TO\_PHYS(\&g\_kernel\_pd))} \\
        从 boot 临时页目录切换到 \texttt{g\_kernel\_pd}。

  \item \textbf{标记就绪}：\texttt{g\_vmm\_ready = 1}

  \item \textbf{初始化 bump allocator}：
        设置 \texttt{g\_next\_vaddr} 为直接映射区上方（对齐到 4\,MiB）。
\end{enumerate}

\codefile{src/kernel/mm/vmm.c}
\begin{lstlisting}[style=cstyle]
void vmm_init(void) {
    printk("[vmm] init: replacing boot PSE with 4KiB page tables...\n");
    memzero_page(&g_kernel_pd);

    // 计算映射范围
    uint32_t kernel_end = (uint32_t)(uintptr_t)__kernel_phys_end;
    uint32_t map_end = kernel_end;
    uint32_t min_map = 16u * 1024u * 1024u;
    if (map_end < min_map) map_end = min_map;
    uint32_t pmm_end = pmm_managed_base()
                      + pmm_total_pages() * VMM_PAGE_SIZE;
    if (pmm_end > map_end) map_end = pmm_end;
    if (map_end > 0x40000000u) map_end = 0x40000000u;

    // 逐页建立双重映射
    for (uint32_t addr = 0; addr < map_end; addr += VMM_PAGE_SIZE) {
        vmm_map_page(addr, addr, PTE_PRESENT | PTE_WRITABLE);
        vmm_map_page(addr + KERNEL_VIRT_OFFSET, addr,
                     PTE_PRESENT | PTE_WRITABLE);
    }

    // 切换到新页目录
    uint32_t pd_phys = VIRT_TO_PHYS(
        (uint32_t)(uintptr_t)&g_kernel_pd);
    vmm_load_page_directory(pd_phys);
    g_vmm_ready = 1;

    // 初始化 bump allocator 水位线
    uint32_t direct_map_top = KERNEL_VIRT_OFFSET + map_end;
    uint32_t align_4m = 0x400000u;
    g_next_vaddr = (direct_map_top + align_4m - 1u)
                 & ~(align_4m - 1u);
}
\end{lstlisting}

\begin{designbox}
\textbf{为什么保留 identity mapping？}

\texttt{vmm\_init} 运行时，内核已经在 high-half（\hex{C0000000}+），按说不需要
低地址映射。但多处仍依赖 identity mapping：
\begin{itemize}
  \item VGA text buffer 在物理 \hex{B8000}，identity map 下可直接写 \texttt{(char*)0xB8000}
  \item Multiboot2 info 指针是物理地址
  \item PMM 内部 bitmap 的物理地址直接当指针用
\end{itemize}
将来 \texttt{vmm\_unmap\_identity()} 会在所有指针转换完成后拆除低地址映射，
把 \texttt{[0, \hex{BFFFFFFF}]} 留给用户态进程。
\end{designbox}

% ============================================================================

